\section{Mixed-integer linear programming formulation}
\label{sec:milp-formulation}

This section records the mathematical model implemented in
\code{circle\_intervals.py}.  It is included so that the computation can be
audited independently of the Python variable names.

\subsection{Variables and interval geometry}

Represent a non-wrapping union of \(N\) intervals by
\[
    A=\bigcup_{i=0}^{N-1}I_i,
    \qquad
    I_i=[x_i,x_i+\alpha_i]\subset[0,1].
\]
The continuous variables satisfy
\[
    0\leq x_i\leq1,
    \qquad
    \alpha_i\geq0,
    \qquad
    x_i+\alpha_i\leq1,
\]
and the ordering constraints are
\[
    x_{i-1}+\alpha_{i-1}\leq x_i
    \qquad(1\leq i<N).
\]
Zero-length intervals are allowed, so an \(N\)-interval model also contains
all configurations with fewer components.

For \(i\leq j\) and any \(\ell\), the real lift of the corresponding sumset
piece is
\[
    I_i+I_j-mI_\ell=[L_{ij\ell},R_{ij\ell}],
\]
where
\begin{align*}
    L_{ij\ell}
      &=x_i+x_j-mx_\ell-m\alpha_\ell,\\
    R_{ij\ell}
      &=x_i+x_j+\alpha_i+\alpha_j-mx_\ell.
\end{align*}

\subsection{Avoiding a point modulo one}

In the standard model, the missing point is fixed to \(t=0\).  In the
translated model, \(t\in[0,1]\) is a continuous variable.  For every triple
\((i,j,\ell)\), an integer variable \(n_{ij\ell}\) selects a unit gap and the
model imposes
\begin{equation}
    n_{ij\ell}+t+\varepsilon\leq L_{ij\ell}
    \leq R_{ij\ell}
    \leq n_{ij\ell}+t+1-\varepsilon.
    \label{eq:lifted-milp-constraint}
\end{equation}
For \(\varepsilon>0\), this places the whole sumset piece a positive distance
from every lift of \(t\).  With \(\varepsilon=0\), equality at an integer
endpoint is allowed.  This is the closed relaxation used to compute a
supremum; it should not be described as strict avoidance for closed intervals.

The objective is
\[
    \text{maximize}\qquad \sum_{i=0}^{N-1}\alpha_i.
\]
The constraints are linear once \(n_{ij\ell}\) is declared integral, so this
is a MILP.

\subsection{Valid cuts and symmetry breaking}

The code can add the following constraints.
\begin{itemize}
    \item The width cut
    \[
        \alpha_i+\alpha_j+m\alpha_\ell\leq1-2\varepsilon
    \]
    is implied by~\eqref{eq:lifted-milp-constraint}.
    \item Ordering implies monotonicity of the lifts: they move right as
    \(i\) or \(j\) increases and left as \(\ell\) increases.  Adding these
    integer inequalities can reduce branch-and-bound work.
    \item In translated mode, translation fixes \(x_0=0\).  One may select a
    shortest interval by imposing \(\alpha_0\leq\alpha_i\) for every \(i\).
    Reflection can then choose \(\alpha_1\leq\alpha_{N-1}\).
    \item In the alternative endpoint convention, reflection is used to add
    \(\alpha_{N-1}\leq\alpha_0\).  It must not be combined blindly with a
    different symmetry convention.
\end{itemize}
These cuts preserve the intended optimum only when their stated symmetry
normalization is valid.  The command-line flags keep them individually
switchable for audit runs.

\subsection{Relaxing selected interactions}

The program may omit an exact triple \((i,j,\ell)\), or every triple supported
on a chosen set of interval indices.  Omitting a triple removes its lift
variable, containment inequalities, width cut, and adjacent monotonicity
links.  The result is a relaxation.  If its certified optimum is still at
most \(1/5\), then the retained constraints suffice for that finite model.  If
its optimum grows, one may only conclude that the omitted interactions matter
to the relaxation.

Subset bounds of the form
\[
    \sum_{i\in S}\alpha_i\leq b
\]
for all subsets \(S\) of a chosen size can also express inductive information,
for example a previously proved bound for every smaller subunion.  This is a
possible bridge from computations at fixed \(N\) to a theoretical induction.

\subsection{Numerical and exact certification}

The implementation uses \code{scipy.optimize.milp}, backed by HiGHS
\cite{SciPy2020,HiGHS2020}.  A trustworthy experiment should archive:
\begin{itemize}
    \item the full command line and git commit;
    \item Python, SciPy, and HiGHS versions;
    \item the serialized variables, constraints, and objective;
    \item primal solution, dual bound, MIP gap, node count, and termination
          status; and
    \item an exact rational check for every claimed construction.
\end{itemize}
A zero requested relative MIP gap still produces a floating-point solver
certificate.  A formal computer-assisted proof would additionally need exact
rational branch verification, interval arithmetic with rigorous error bounds,
or an exact solver whose certificate can be checked by a small independent
program.

\subsection{MILP around the equally spaced base}

The companion script \code{equally\_spaced\_tiny\_interval\_lp.py} fixes the
base parameters \(q,\delta,a\) from
Section~\ref{sec:equally-spaced-construction}, adds variable extra intervals,
and imposes the same lifted constraints.  Pairwise binary variables enforce
non-overlap because the extra intervals are not assigned to gaps in advance.
Its role is local discovery: if the solver finds a positive extra length, the
next step is to recover rational endpoints and produce a finite exact
certificate.
